% =============================================================================
%  Formatting.tex — shared macros, packages, and styles for the report
% =============================================================================

\usepackage{listings}
\usepackage{amssymb}
\usepackage{xcolor}
\usepackage{booktabs}
\usepackage{array}
\usepackage{longtable}
\usepackage{multirow}
\usepackage{tcolorbox}
\tcbuselibrary{skins,breakable}
\usepackage{enumitem}
\usepackage{siunitx}
\usepackage{tikz}
\usetikzlibrary{positioning, arrows.meta, shapes.geometric, calc, fit, backgrounds}

% ---------------------------------------------------------------------------
% Table breathing room — prevents wrapped multi-line cells in narrow p{}
% columns from visually crowding/overlapping adjacent rows.
% ---------------------------------------------------------------------------
\renewcommand{\arraystretch}{1.25}

% Ragged-right paragraph column (avoids the huge stretched inter-word gaps
% that full justification produces in narrow table columns).
\newcolumntype{L}[1]{>{\raggedright\arraybackslash}p{#1}}

% ---------------------------------------------------------------------------
% Colour palette (mission-control theme, used sparingly in print)
% ---------------------------------------------------------------------------
\definecolor{codebg}{RGB}{245,245,248}
\definecolor{codeframe}{RGB}{200,200,210}
\definecolor{accentblue}{RGB}{59,116,228}
\definecolor{alertred}{RGB}{192,20,46}
\definecolor{caseboxbg}{RGB}{240,244,250}

% ---------------------------------------------------------------------------
% Code listings style (Python / bash / JSON snippets throughout the report)
% ---------------------------------------------------------------------------
\lstdefinestyle{codestyle}{
  backgroundcolor=\color{codebg},
  rulecolor=\color{codeframe},
  frame=single,
  basicstyle=\ttfamily\footnotesize,
  keywordstyle=\color{accentblue}\bfseries,
  commentstyle=\color{gray}\itshape,
  stringstyle=\color{teal},
  numbers=left,
  numberstyle=\tiny\color{gray},
  numbersep=6pt,
  breaklines=true,
  breakatwhitespace=false,
  showstringspaces=false,
  tabsize=2,
  captionpos=b,
  xleftmargin=1.2em,
  framexleftmargin=1em,
}
% ---------------------------------------------------------------------------
% listings ne fournit pas JavaScript. Defini ici pour l'annexe D (logique de
% selection du transport WebSocket du cockpit), plutot que de rabattre sur
% `Java`, dont les mots-cles ne correspondent pas.
% ---------------------------------------------------------------------------
\lstdefinelanguage{JavaScript}{
  keywords={const, let, var, function, return, if, else, for, while, new,
            typeof, instanceof, this, null, undefined, true, false, async,
            await, class, extends, import, export, default, try, catch},
  ndkeywords={window, document, console, location, JSON, Math, Promise},
  sensitive=true,
  comment=[l]{//},
  morecomment=[s]{/*}{*/},
  morestring=[b]',
  morestring=[b]",
  morestring=[b]`,
}

\lstset{style=codestyle}

% ---------------------------------------------------------------------------
% Case-study box (used extensively in Chapter 6 — debugging case studies)
% ---------------------------------------------------------------------------
\newtcolorbox{casestudy}[1]{
  breakable, enhanced,
  colback=caseboxbg, colframe=accentblue,
  fonttitle=\bfseries, title={#1},
  boxrule=0.6pt, arc=2pt,
  left=6pt, right=6pt, top=4pt, bottom=4pt,
}

% ---------------------------------------------------------------------------
% Trap box (Appendix D) --- for the failure modes that cost days and whose
% symptom points at the wrong subsystem. Deliberately red and deliberately
% distinct from \texttt{casestudy}: a case study narrates a resolved
% investigation, a trap warns a reader who is about to lose an afternoon.
% ---------------------------------------------------------------------------
\newtcolorbox{trapbox}[1]{
  breakable, enhanced,
  colback=alertred!4, colframe=alertred,
  fonttitle=\bfseries, coltitle=white, title={#1},
  boxrule=0.6pt, arc=2pt,
  left=6pt, right=6pt, top=4pt, bottom=4pt,
}

% ---------------------------------------------------------------------------
% Equation-block macro — for formulas that were laid out as single-cell
% tables in the source document; rendered here as proper LaTeX math.
% ---------------------------------------------------------------------------
\newenvironment{eqblock}
  {\begin{center}\begin{minipage}{0.92\linewidth}}
  {\end{minipage}\end{center}}

% ---------------------------------------------------------------------------
% Short helper macros used throughout
% ---------------------------------------------------------------------------
% \code allows LaTeX to break long identifiers (e.g. HEARTBEAT_STALE_S) after
% each underscore when they don't fit a narrow table column — without this,
% such unbreakable "words" overflow their column and visually collide with
% the next one. Locally redefining \_ (rather than using seqsplit, which
% chokes on the \_ control sequence) keeps every other use of \_ in the
% document untouched.
\newcommand{\breakableunderscore}{\textunderscore\penalty0\hskip0pt\relax}
% \DeclareRobustCommand (not \newcommand) so \code survives moving arguments
% such as \caption{...} and \section{...}, where a fragile \renewcommand
% assignment would otherwise raise "Argument ... has an extra }".
% \texttt (an NFSS text command, not a bare declaration like \ttfamily) is
% used as the outer wrapper so \code also remains valid inside math mode.
\DeclareRobustCommand{\code}[1]{\texttt{\renewcommand{\_}{\breakableunderscore}#1}}
\DeclareRobustCommand{\topic}[1]{\texttt{\renewcommand{\_}{\breakableunderscore}#1}}

% \todo{...} — inline placeholder for a value the operator fills in after a
% field run (e.g. a measured number not yet available at write time).
% Renders as a bracketed, coloured tag so it cannot be missed in the
% compiled PDF, and is easy to grep for in the source ("\todo{").
\newcommand{\todo}[1]{\textcolor{alertred}{\textbf{[TODO: #1]}}}
